%%%%%%%%%%%%%%%%%%%%%%%%%%%%%%%%%%%%%%%%%%%%%%%%%%%%%%%%%%%%%%%%%%%%%%%%%%%
% Función: capítulo de ejemplo de trabajos previos de la tesis por compendio.
%
% Uso: copia de restauración; normalmente no debes modificarla.
% Tras restaurarla, modifica la copia activa para adaptarla a tu trabajo.
%%%%%%%%%%%%%%%%%%%%%%%%%%%%%%%%%%%%%%%%%%%%%%%%%%%%%%%%%%%%%%%%%%%%%%%%%%%

\ifthenelse{\equal{\myLanguage}{english}}
{
  \chapter{Previous Work}
}
{
  \chapter{Trabajo previo}
}
\label{cha:compendium-previous-work}

\ifthenelse{\equal{\myLanguage}{english}}
{
  A literature review can distinguish a specific visual localisation method~\cite{Sanabria2026Landmark}, a functional-assessment proposal~\cite{nieva2024proposal} and a detection and localisation application using distributed acoustic sensing~\cite{ramirez2026jstars}. The review of machine-learning methods for pipeline surveillance~\cite{tejedor2017machine} also provides an example of a survey rather than an individual method paper. These examples illustrate ways to organise sources, but they do not establish that all the works address the same research question.
}
{
  Una revisión bibliográfica puede distinguir un método concreto de localización visual~\cite{Sanabria2026Landmark}, una propuesta de evaluación funcional~\cite{nieva2024proposal} y una aplicación de detección y localización mediante sensado acústico distribuido~\cite{ramirez2026jstars}. La revisión de métodos de aprendizaje automático para vigilancia de tuberías~\cite{tejedor2017machine} ofrece además un ejemplo de revisión del estado del arte, en lugar de un artículo sobre un método individual. Estos ejemplos ilustran formas de organizar las fuentes, pero no implican que todos los trabajos aborden la misma pregunta de investigación.
}

\ifthenelse{\equal{\myLanguage}{english}}
{
  \section{Scientific Background}
  \subsection{Foundational approaches}
}
{
  \section{Antecedentes científicos}
  \subsection{Enfoques fundamentales}
}

\ifthenelse{\equal{\myLanguage}{english}}
{
  Review the scientific and technical background of the thesis here. Organize the discussion around the research questions rather than presenting an isolated list of publications, and identify the limitations that motivate the doctoral work.
}
{
  Revisa aquí los antecedentes científicos y técnicos de la tesis. Organiza la discusión alrededor de las preguntas de investigación en lugar de presentar una lista aislada de publicaciones, e identifica las limitaciones que motivan el trabajo doctoral.
}

\ifthenelse{\equal{\myLanguage}{english}}
{
  \subsection{Recent developments}

  Summarize the approaches most directly related to the thesis and compare their assumptions, strengths and limitations.

  \section{Research Gap}

  Conclude the review by identifying the unresolved questions addressed by the objectives of this thesis.
}
{
  \subsection{Desarrollos recientes}

  Resume los enfoques más relacionados con la tesis y compara sus hipótesis, fortalezas y limitaciones.

  \section{Oportunidad de investigación}

  Concluye la revisión identificando las cuestiones no resueltas que abordan los objetivos de esta tesis.
}

%%% Local Variables:
%%% TeX-master: "../../book"
%%% End:
